\section{L'impasse ergonomique}
Si l'infrastructure technique de CujasNum pose d'importants défis en termes de sécurité, la limite se révèle dans son fonctionnement quotidien. La chaîne de numérisation actuelle de la bibliothèque numérique illustre parfaitement l'approche décrite par Elina Leblanc dans sa thèse la qualifiant d'une \enquote{approche purement descendante(top-down), axée sur les contraintes du système plutôt que sur l'activité des agents}.

Elina Leblanc s'appuie sur les critères d'Aaron Schmidt, l'évaluation d'un service de bibliothèque numérique présente \enquote{un guide pour la CCU des bibliothèques analogiques [...]sous un angle utilisateur}. Il mentionne ainsi l'utilité d'un service d'une bibliothèque numérique en disant que celui-ci doit être \enquote{utilisable et désirable}. Aaron Schmidt l'explique ainsi : \enquote{Un service utile doit répondre à un besoin et avoir une importance aux yeux des utilisateurs pour qu'ils ne l'utilisent pas de manière ponctuelle, mais sur la longue durée. Un service utilisable garantit aux utilisateurs un accès aisé et rapide aux données. Enfin, la désirabilité crée un lien entre l'utilisateur et un service, en donnant envie au premier d'utiliser le deuxième et en suscitant un besoin}. Cela garantit à l'utilisateur une expérience optimale lorsqu'ils utilisent la bibliothèque numérique d'une institution patrimoniale. Or, à la bibliothèque Cujas, l'interface n'est pas adapté à la manière dont l'humain utlise les sites web aujourd'hui. L'utilisateur doit s'adapter lui aux contraintes rigides de l'interface et à la façon dont lui sont présentes les services. 

Ce manque d'ergonomie se traduit par un décalage dans la production des ouvrages numériques. Si le contrôle qualité d'un ouvrage requiert environ \enquote{quinze minutes en moyenne si tout va bien}, la moindre anomalie structurelle ou typographique fait basculer le traitement dans une autre dimension nécessitant au minimum \enquote{trois jours} de manipulations afin de corriger l'erreur détectée et ensuite pouvoir remettre en ligne le document. Cette manière de travailler est devenue un gouffre chronophage et la rigidité qu'implique ce traitement est devenue un carcan dont la bibliothèque Cujas ne parvient pas à se défaire.

\subsection{La pénibilité des traitements manuels et de l'encodage XML}
Sous CujasNum, la processus de production ne bénéficie d'aucune automatisation globale, s'apparentant plutôt à un \enquote{bricolage permanent} destiné à préserver l'autonomie des équipes tout en maintenant une indispensable supervision humaine. Le premier goulot d'étranglement de cette chaîne réside dans l'obligation de manipuler plusieurs formats et de les maintenir simultanément pour une même donnée structurelle. C'est précisément le cas pour leur table des matières d'un ouvrage, qui sont retranscrite manuellement afin de reconstituer intellectuellement l'ouvrage, un élément d'ailleurs essentiel pour permettre à l'utilisateur de naviguer à l'intérieur d'un document numérisé. Cette table des matières est saisie dans un fichier Excel, avant que celui-ci ne soit utilisé pour l'exporter au format CSV, initiant ainis une chaîne de conversions manuelles fastidieuses.

À cela s'ajoute l'obligation de respecter un protocole d'édition extrêmement rigide, tel que l'obligation de nommer le premier onglet du tableau Excel \enquote{Feuille 1} ou de vérifier, livre en main, la stricte concordance entre les images numériques et la pagination réelle de l'ouvrage. Le coût cognitif le plus important apparaît lors de la validation du code XML. Afin d'éviter un rejet par le CINES, les agents doivent traquer manuellement les moindres anomalies typographiques : remplacer les guillements français par des guillements anglais pour ne pas perturber le parseur XML, ou encoder de force les esperluettes sous la forme de l'entité \enquote{\&amp};. Ils doivent veiller à ce que l'arborescence ne comporte aucun saut hiérarchique descendant incohérent -- par exemple passer \enquote{directement au d'un niveau hiérarchique 1 (Chapitre) à un niveau 3 (sous-section) sans déclarer formellement le niveau 2 intermédiaire}. 

Par ailleurs, le techinique en charge du contrôle doit en plus de s'assurer de la conformité pour le CINES, \enquote{vérifier qu'il y a bien une image pour chaque page} ainsi que le nommage du fichier image nommé avec \enquote{le code-barres suivi du numéro d'image et de l'extension}. Ensuite, il doit regarder si les liens commençant par \enquote{http} ont été remplacé par \enquote{https}. Toutes ces micro-corrections, indispensables pour empêcher le blocage lors de l'affichage ou du versement des fichiers reposent sur des interventions humaines précises et mettent en avant cette idée d'un \enquote{bricolage manuel permanent} s'assurant de la conformité des notices. L'infrastructure technique ne propose pour le moment aucune interface d'aide à la saisie ou de correction automatique afin d'éviter aux équipes ce travail méticuleux et chronophage. 

Ces micro-intervention créent une défaillance ergonomique entraînant une asymétrie temporelle dans les différentes étapes de la chaîne de numérisation. En temps normal, le contrôle qualité d'un ouvrage dure environ \enquote{quinze minutes en moyenne si tout va bien}. En revanche, lorsque les scripts détectent des erreurs cela peut prendre proportion inadapté contenue de l'erreur initiale. Habituellement, la correction d'une erreur simple, comme un guillemet, devrait être proportionnelle à sa gravité, or, à Cujas, cela n'est pas le cas. Actuellement l'interface n'offre aucune possibilité de faire une modification rapide \textit{in situ} en évitant d'avoir à relancer toute la procédure pour une coquille. En effet, l'apparition d'une métadonnée erronnée ou d'un bug oblige en l'état, l'agent à \enquote{extraire le document, à corriger manuellement les fichiers METS ou Dublin Core sous Oxygen, puis à ré-uploader le lot via un protocole SFTP sécurisé (WinSCP) dont le traitement par lot n'est planifié que durant la nuit}. La rectification d'une coquille ou d'un lien brisé peut nécessiter jusqu'à trois jours de traitement minimum avant d'être visible en ligne pour l'usager.  

\subsection{La rigidité d'XTF face aux formats non monographiques}
À l'origine l'architecture d'XTF a été pensée exclusivement pour un seul type de document -- la monographie -- se révèle incapabe de s'adapter à d'autres types de documents avec des arborescences complexes, comme les périodiques. XTF repose sur une arborescence dite \enquote{à plat}. C'est la raison pour laquelle si on ajoute d'autres types de documents l'arobrescence ne sera pas respecté, constituant un sérieux problème pour les équipes de Cujas qui souhaite ajouter d'autres collections sur leur bibliothèque numérique. 

XTF possède de nombreuses contraintes techniques, l'une d'elles réside dans l'obligation d'avoir le format PDF afin de mettre le document en ligne. Le système a été conçu comme \enquote{un simple répertoire PDF} calqué sur la structure physique d'un livre (volume et pages). Cela rend donc difficile l'ajout de tout autre types de documents ne respectant pas ce principe. Pour rendre la consultation des images plus agréables, une visionneuse JPEG a été implémenté en 2013 pour permettre à l'usager de feuilleter le document de façon dynamique. 

Contrairement à une monographie, un périodique exige une architecture d'accès d'au moins de trois niveau, \enquote{c'est-à-dire le titre de la revue, suivie de l'année de publication et du numéro du volume ou du fascicule, pour que la consultation soit pertinente}. Etant incapble de reproduire cette structure, c'est l'une des raisons qui pousse la bibliothèque Cujas à mettre de côté ses collections ne pouvant les intégrer sans \enquote{dénaturer} le type de document. Aujourd'hui, on observe sur certaines bibliothèques numériques possédant des périodiques que la granularité peut aller jusqu'à l'article même. XTF s'arrête initialement au volume et ne permet pas une granularité aussi précise, allant donc à l'encontre des pratiques de recherche actuelle. Faute de pouvoir diffuser ces revues dans sa bibliothèque numérique, Cujas a fait \enquote{le choix de faire numériser ces volumes via un partenariat avec Persée}. Pour signaler ces numérisations, la bibliothèque se contente uniquement de créer des notices de redirection indiquant aux chercheurs l'existance de ces documents où est-ce qu'il peut les trouver, afin de pouvoir les consulter. 

La bibliothèqe Cujas rencontre le même problème avec son fonds iconographiques qui ne possédant pas de structure intellectuelle répliquable sous la forme d'une table des matières. Si les équipes de Cujas souhaitent mettre en ligne ce fonds. Cela peut être effectuée \enquote{de manière détournée en créant de PDF}. Ils doivent alors créer \enquote{artificiellement une table des matières et une table de concordance}, afin de pouvoir mettre en ligne ces documents. Cependant, c'est au niveau de l'archivage pérenne que le blocage risque de se faire suite à ce type de bricolage. En effet, cela obligerait la bibliothèque Cujas à \enquote{passer par une redéfinition complexe du plan d'archivage avec le CINES}. 

Un autre problème réside dans l'indexation utilisée par la bibliothèque numérique actuelle, l'indexation Hauville. Celle-ci ne permet pas le filtrage via différents types de documents; en effet elle a été structuré qu'autour de la monographie ne permettant pas de croiser les informations. L'impossibilité d'ajouter d'autres types de documents met en lumière l'invisibilisation auquel ils se trouveraient confrontés lors d'une recherche plein-texte. Ces documents ne pourraient ressortir de manière évidente ne pouvant posséder d'OCR. Les résultats apparaitront donc de façon très restreinte et dans les dernières pages problablement. L'obsolescence du logiciel XTF entraîne donc \textit{de facto} un filtrage empêchant la bibliothèque Cujas de pouvoir mettre en avant la richesse et la diversité de ses collections. Il lui est impossible de valoriser ses documents afin de justifier les choix effectués ou de mettre en avant une collection ou un document dans un article ou dans une exposition virtuelle. 
\medskip

Le bilan de CujasNum met en avant les points forts du projet ainsi que les points faibles au moment de la réalisation de celui-ci, mais également ceux ayant émergé par la suite avec l'évolution des technoloqies numériques et des pratiques du Web. D'abord, pionnière dans le monde des bibliothèques numériques, la bibliothèque Cujas s'est trouvé progressivement isolée parvenant difficilement à suivre l'évolution rapide liée à l'émergence de l'internet entraînant les pratiques documentaires et les pratiques d'Internet dans un nouvel âge. Face à ce constat, \enquote{l'ajout d'une surcouche WordPress ne s'aurait être la solution comme cela avait été envisagé dans un premier temps}. Pour répondre aux attentes contemporaines des chercheurs et pour s'adapter aux mutations nationales des systèmes de gestion des bibliothèques (SGB) impose une rupture dans les pratiques actuelles. Cela nécessite de réfléchir à une interface plus souple, acceptant différents types de documents se basant sur les pratiques actuelles de la recherche, répondant aux nouveaux enjeux des utilisateurs. La bibliothèque numérique doit permettre à sa bibliothèque d'utiliser les standards du Web de données afin de rendre sa plateforme davantage interopérable et prête à s'intégrer dans de grands réseaux scientifiques nationaux.